\section{ノートシリーズとの関係}
\label{sec:note}

本論文のsourceは、沈殿、流れ、地形の語彙を用いる関連ノートシリーズ（L1層）に言及している。
正確なノート成果物、題名、版は未解決である。
表は、記録された語彙とモデル構造についてのCT自身の教育的比較を保持する。
独立に同定されたノートの検証済みの再構成でも、型を明示した写像でも、検証のための証拠でもない。
認知的・現象的な項目は解釈の提案であり、表示した対応によって確立された同一性ではない。

\begin{center}
\renewcommand{\arraystretch}{1.5}
\begin{tabular}{p{0.40\textwidth}p{0.52\textwidth}}
\toprule
\textbf{ここに記録されたノート語彙（J：日本語）} &
\textbf{提案するモデル対応} \\
\midrule
「ログの沈殿」 & $\sigma_{\mathrm{drive}}$を介した$\Phi$の深化 \\
「安定点」 & 局所的安定領域：$\nabla\Phi\approx 0$ \\
「概念のくぼみ」 & $\Phi$のポテンシャル井戸 \\
「全体勾配」 & 情報ポテンシャル$\Phi = K*I$ \\
「通りやすさ」 & $-\Phi(x,\tau)$（通りやすさが高い$=$$\Phi$が低い） \\
「谷」 & $\Phi$の深い井戸 \\
「壁」 & $\Phi$が高い尾根領域 \\
「理解した感覚」 & 通過コスト$\int|\nabla\Phi|dx$の低下 \\
「思考の癖」 & $\JI$のリミットサイクル軌道 \\
「GPU的な装置」 & $\JI$に対する$\Phi$の閾値未満の作用 \\
「言語粒子」 & $\nabla\Phi$へ結合する離散トークン$\ell_i$ \\
「圧縮」 & 量子化作用素$\mathcal{Q}$ \\
「社会的共有地形」 & エージェント間での社会的$\Phistar$の維持の提案 \\
\bottomrule
\end{tabular}
\end{center}

\subsection{この比較の限界}
ここに記録された記述は、未同定のノート成果物に検証済みの立場を帰属させることなく、非専門的な直観として用いる。
CTの立場では、それらは場理論的量の\emph{解釈}であり、定義ではない。

特に、
\begin{itemize}
\item 記録された意味の語彙は、反復的な通過とポテンシャル井戸の深さと比較する。
内在的な意味対象や、較正された意味尺度と比較するものではない。
\item 記録された「理解した感覚」の記述には、\eqref{eq:understanding}に関連する容易さの変化との対応を提案する。
局所的定義は、この現象的対応の独立した証拠ではない。
\item ここでConsciousnessのノートに帰属させているCPU/GPUの区別を、高速レーン／低速レーンの分解と比較する。
Consciousness Theoryには現在repository内のsourceがあるが、ノートと論文の正確な版の関係は未解決である。
これは構造的比較であり、共通量、型を明示した翻訳、実証された機構の移転ではない。
\end{itemize}

\subsection{L3 → L1の方向}
本論文（L3）はノートシリーズ（L1）を改訂したり置き換えたりしない。
ここでの寄与は解釈の整理であり、記録された直観を比較するためのモデル語彙であって、それらの理論的検証ではない。
正確な出典の帰属は未解決である。

地形、流れ、沈殿を場理論的対象と比較する。
理解、意味、思考はその語彙の別々の解釈的な読みであり、形式的な指示対象を割り当てることによって独立に説明されるものではない。
二つの説明層は、一方が他方の検証証拠となることなしに、相補的でありうる。
